\subsection{An explicit reciprocity obstruction to unrestricted cyclic repair}

\begin{statement}{A.110}{Refuted}
{Unrestricted prescribed local classes with cyclic auxiliary ramification}
The following unrestricted strengthening of (A.109.4) is false:
one may prescribe arbitrary local norm-one classes at a finite set
and require all newly ramified primes to split the coefficient field
and have cyclic local image.  It is still false when, at each new
split prime, the allowed condition is enlarged to the exact
scalar-shifted graph $c=zt+w\mathbf1$ of (A.109.8).

An explicit example uses
\[
 a=3,\quad N=9,\quad r=1,\quad
 K=\Q(3^{1/9}),\quad L=\Q(3^{1/9},\mu_9),\quad
                 q=271,\quad \eta=6.
 \tag{A.110.1}
\]
Index the embeddings of $K$ by $i\in\Z/9\Z$ and use the module
$M=I_1\otimes\mu_3$ of A.109.
At $q$ prescribe
\[
 \gamma_q=(e_0-e_1)\,\kappa_q(q)
                   +(e_3-e_2)\,\kappa_q(6),
                   \qquad \gamma_2=\gamma_3=0.
 \tag{A.110.2}
\]
Every displayed local class comes from $T(\Q_v)/T(\Q_v)^3$.
They have global preimages under finite-set localization,
and even global norm-one representatives by the finite norm-torus
weak approximation used in A.101.
No such global class has the required cyclic or scalar-shifted
graph condition at every new ramified prime splitting $L$,
while being unramified at every other prime outside
$\{2,3,271\}$.

This conclusion does not refute the canonical ordinary norm-chain
repair.  In particular the quadratic invariant used here is
surjective on the actual local norm-one group $T(\Q_3)/T(\Q_3)^3$.
Leaving that local multiplier free can therefore compensate for
any value of this particular invariant.
\end{statement}

\begin{proof}
Define $f:\Z/9\Z\to\F_3$ and a form on the permutation module by
\[
 f(d)=\begin{cases}d/3\pmod3&3\mid d,\\0&3\nmid d,\end{cases}
 \qquad
 \mathcal B(v,w)=\sum_{i,j}v_iw_j f(j-i).
 \tag{A.110.3}
\]
Here $d/3\pmod3$ is well-defined in the first case.
We have $f(-d)=-f(d)$ and
$f(ud)=(u\bmod3)f(d)$ for every unit $u\bmod9$.
Translations preserve differences.  The two cyclotomic twists
on $M$ multiply to one modulo three, and hence restriction gives
the actual equivariant alternating map
\[
             \mathcal B:\mathop{\bigwedge}\nolimits^2M\longrightarrow\mu_3.
 \tag{A.110.4}
\]
Moreover $\sum_d f(d)=0$, so $\mathcal B(\mathbf1,v)=0$ for
every $v$, and
\[
 \mathcal B(e_0-e_1,e_3-e_2)=1.
 \tag{A.110.5}
\]
For a global class $\gamma\in H^1(\Q,M)$ define the actual
Brauer class
\[
       \mathcal Q(\gamma)=\tfrac12\mathcal B_*(\gamma\smile\gamma)
                  \in H^2(\Q,\mu_3)=\operatorname{Br}(\Q)[3].
 \tag{A.110.6}
\]
This is independent of the cocycle representative by the
functoriality of cup product.  The factor $1/2$ is the inverse
of two in $\F_3$.

At a split prime with local coordinates
$\gamma=t\kappa(q)+c\kappa(\eta)$, graded commutativity of cup
product and alternation of $\mathcal B$ give
\[
       \mathcal Q(\gamma_q)
             =\mathcal B(t,c)\,\kappa(q)\smile\kappa(\eta).
 \tag{A.110.7}
\]
This also follows from the literal normalized cocycle
$\tfrac12\mathcal B(\gamma(g),g\gamma(h))$.
For a cyclic local image, $t,c$ are dependent over $\F_3$;
for a scalar-shifted graph, $c=zt+w\mathbf1$.
In either case the coefficient in (A.110.7) is zero.

The prime in (A.110.1) satisfies
\[
       271\equiv1\pmod9,\qquad 38^9\equiv3\pmod{271},
                    \qquad 6^{90}\equiv242\not\equiv1\pmod{271}.
 \tag{A.110.8}
\]
Hensel lifting shows that it splits $L$.
The unit six generates the cubic unit quotient, so the local
class $\kappa(q)\smile\kappa(6)$ has nonzero Brauer invariant.
For example, $\Q_q(6^{1/3})/\Q_q$ is unramified of degree three,
whose norm valuations are divisible by three; its Hilbert pairing
with $q$ is therefore nonzero.
Equations (A.110.2), (A.110.5), and (A.110.7) give a nonzero
invariant at $271$.

At two and three the prescribed class is zero.
The coefficient field $L$ is unramified away from three;
at every other unramified place the cup class is inflated from
the residue-field Galois group, whose finite three-primary
cohomology in degree two is zero.
The real place has no three-primary cohomology in positive degrees.
Every allowed new ramified prime gives zero by (A.110.7).
Thus a hypothetical global class satisfying the strengthening
would give a Brauer class with precisely one nonzero local
invariant.  This contradicts global reciprocity, equivalently
the final map in the finite duality sequence
\cite[I, Theorem 4.10(c)]{BIR_MilneADT}.

The asserted local norm-one interpretation is literal:
at $271$, the tuple with coordinates $q$ at index zero,
$q^{-1}$ at index one, $6$ at index three, $6^{-1}$ at
index two, and one elsewhere has product one.
It represents (A.110.2); at two and three use the identity.
Finite-set approximation for $T$ supplies an actual global
norm-one element with these power classes, as stated.
The proof establishes the impossibility of the added
ramification condition for that local prescription.

\emph{The exact remaining freedom at three.}
Write the translation group algebra as
$\F_3[C_9]=\F_3[Z]/(Z^9)$ with $Z=x-1$.
The basis $Z,Z^2,\ldots,Z^8$ identifies $I_1$.
In this basis the form (A.110.3) has its only nonzero block
\[
 \begin{pmatrix}
 0&0&2&2\\0&0&1&0\\1&2&0&0\\1&0&0&0
 \end{pmatrix},
 \qquad
 \operatorname{rad}(\mathcal B)=
                    Z^5\F_3[Z]/(Z^9)\otimes\mu_3=:R_0.
 \tag{A.110.9}
\]
The block is invertible.  Let $V=M/R_0$; the induced alternating
pairing $V\times V\to\mu_3$ is nondegenerate.
The local group at three maps onto the full holomorph $G$,
as in A.105.
The translation-fixed line in $V$ is represented by $Z^4$;
the involution has its cyclotomic twist $-1$ on that line.
Hence $H^0(\Q_3,V)=0$.
The translation coinvariant of $R_0$ is represented by $Z^5$;
the same involution acts there by $(-1)^5(-1)=1$.
Every map from $R_0$ to $\mu_3$ must factor through that
coinvariant, so $H^0(\Q_3,R_0^\vee)=0$.
Local duality \cite[I, Corollary 2.3]{BIR_MilneADT}
gives $H^2(\Q_3,R_0)=0$ and therefore
\[
       H^1(\Q_3,M)\longrightarrow H^1(\Q_3,V)
                    \quad\text{is onto}.
 \tag{A.110.10}
\]
Since $V\simeq V^\vee$, local duality and the local Euler
characteristic \cite[I, Theorem 2.8]{BIR_MilneADT}
give $\dim_{\F_3}H^1(\Q_3,V)=4$ and a
nondegenerate symmetric pairing on this space, obtained by
cup product followed by $\mathcal B$ and the local invariant.
Its half-diagonal is a nondegenerate quadratic form.
A nondegenerate quadratic form of dimension at least two over
$\F_3$ represents every element: diagonalize two variables;
opposite coefficients already give both nonzero values, and
equal coefficients give their two different sums.
Thus (A.110.10) proves surjectivity of $\mathcal Q$ on
$H^1(\Q_3,M)$.

Finally $N_{\Q_3(3^{1/9})/\Q_3}$ is onto by composing the two
non-Galois cubic norm maps.  The first cubic has $S_3$ closure
over $\Q_3$; the second is the map of A.105 with $Q=1$.
In both cases the index-two closure norm subgroup together
with the scalar cubes proves surjectivity, exactly as there.
Consequently $H^1(\Q_3,T)=0$, and the local Kummer sequence
identifies $T(\Q_3)/T(\Q_3)^3$ with $H^1(\Q_3,M)$.
This proves the stated surjectivity using actual local
norm-one elements.  It does not prescribe their values at
other places or construct a compatible global norm chain.
\end{proof}

For the exact multiplier $u(h_{\rm bad})$ of A.106, the local
vectors at $321193$ are $t=e_0-e_1$ and $c=e_3-e_4$.
The same map gives $\mathcal B(t,c)=2$.
Thus A.106 has a nonzero contribution to this very Brauer class;
reciprocity forces another contribution for its actual global
multiplier.  Its place is not identified here.

There is also an exact relative limitation.  The norm map from
$K=\Q(3^{1/9})$ to $K_1=\Q(3^{1/3})$ sends geometric
coordinates to their products over indices congruent modulo three.
On finite torsion its map is therefore the block sum
$P_9\to P_3$.  The relative norm-one module is
\[
 M_{\rm rel}=\ker(P_9\longrightarrow P_3)\otimes\mu_3
       =Z^3\F_3[Z]/(Z^9)\otimes\mu_3.
 \tag{A.110.11}
\]
The form $\mathcal B$ vanishes identically on this submodule.
Indeed on a block write $i=h+3i'$, $j=h+3j'$; its contribution
is
\[
 \Bigl(\sum_{i'}v_{h+3i'}\Bigr)
       \Bigl(\sum_{j'}j'w_{h+3j'}\Bigr)
 -\Bigl(\sum_{i'}i'v_{h+3i'}\Bigr)
       \Bigl(\sum_{j'}w_{h+3j'}\Bigr)=0,
 \tag{A.110.12}
\]
and different blocks have zero pairing.
Thus this specific invariant imposes no condition on a change
whose lower norm class is preserved through this relative module.
The A.106 vector $e_0-e_1$ has nonzero lower block sums.
The comparison (A.110.11) retains the lower-norm map and is
independent of the separate freedom at three.

The shared retained certificate
\texttt{norm\_torus\_selmer\_certificate.py}
checks all fifty-four arithmetic equivariance identities,
59,049 scalar-shifted graph isotropy identities, 729 literal
normalized two-cocycle identities, and the explicit splitting
and cubic residue data at $271$.
